This paper quantifies the total factor productivity (TFP) gain required to offset the short-run output loss associated with alternative weekly working-hour ceilings in Brazil. We combine moments from the 2024Q4 PNAD Contínua with a static partial-equilibrium model in which formal and informal labor are imperfect substitutes, hourly efficiency depends on working hours, and formalization wedges govern the composition of employment. To represent reductions from 44 to 40 and 36 hours, we initially cap formal workers' usual hours at 44 as a baseline measurement assumption. The 40-hour alternative requires $1.53$--$1.57\%$ higher TFP. Moving to 36 hours requires $6.52$--$8.38\%$, depending on the efficiency assumption below 40 hours; informality rises by $2.28$--$2.92$ percentage points in the same counterfactual. Decompositions show that fewer physical formal hours account for the largest negative component of the output change. Comparison with Brazil's historical TFP path and welfare exercises suggests an important quantitative difference between moderate reforms and the 36-hour ceiling. Magnitudes depend on fatigue and baseline hours: retaining the observed usual-hours tail above 44 reduces the requirement at 40 hours to approximately $0.07\%$. The results motivate the assessment of gradual paths and transition instruments, without providing a causal estimate of the effects of a specific legislative proposal.
